%!TEX program = lualatex
% SPDX-License-Identifier: CC-BY-SA-4.0
% Requires: https://github.com/jemmybutton/fiziko (GPL-3.0)
% Set MPINPUTS to the directory containing fiziko.mp before running lualatex.
\documentclass[11pt]{article}

\usepackage{iftex}
\ifLuaTeX\else
  \PackageError{vintage-latex}{This example requires LuaLaTeX}
    {Run lualatex, not pdflatex or xelatex.}
\fi

\usepackage{fontspec}
\usepackage{unicode-math}
\usepackage{xcolor}
\usepackage{geometry}
\usepackage[tracking=true,protrusion=false,expansion=false]{microtype}
\usepackage{luamplib}

\setmainfont{EB Garamond}[Numbers={OldStyle,Proportional}]
\setmathfont{Garamond-Math.otf}
\definecolor{ragpaper}{HTML}{F3E9D6}
\definecolor{printerink}{HTML}{25201B}
\pagecolor{ragpaper}
\color{printerink}
\geometry{a5paper,landscape,top=16mm,bottom=16mm,left=18mm,right=18mm}
\pagestyle{empty}
\setlength{\parindent}{0pt}

\mplibforcehmode
\everymplib{
  input fiziko.mp;
  randomseed := 11;
  defineMinStrokeWidth(0.2pt);
  defineLightDirection(-1/8pi, 1/8pi);
  color inkcolor; inkcolor := (0.145, 0.125, 0.106);
  drawoptions(withcolor inkcolor);
  beginfig(1);
}
\everyendmplib{endfig;}

\newcommand{\knotcell}[2]{% figure code, caption
  \begin{minipage}[t]{.31\linewidth}
    \centering
    \begin{mplibcode}
      #1
    \end{mplibcode}\par
    \vspace{.5em}
    {\small\itshape #2}
  \end{minipage}%
}

\begin{document}

\begin{center}
  {\small\textls[120]{\textsc{Plate IX}}}\par
  \vspace{.3em}
  {\Large Of Knots}\par
  \vspace{.2em}
  {\itshape three closed cords, drawn as the sailor sees them}
\end{center}

\vspace{1.4em}

% Each cord is a closed curve given by a formula, sampled every few degrees.
% fiziko finds the crossings itself; the string of ones tells it which strand
% passes over at each crossing in turn, alternating as in a real knot.
\knotcell{
  path cord;
  cord := (for t = 0 step 12 until 348:
    (sind(t) + 2sind(2t), cosd(t) - 2cosd(2t)) scaled 0.5cm .. endfor cycle);
  addStrandToKnot (trefoil) (cord, 0.24cm, "l", "1,-1,1,-1,1,-1");
  draw knotFromStrands (trefoil);
}{Fig.\ 1. The trefoil, of three crossings.}
\hfill
\knotcell{
  path cord;
  cord := (for t = 0 step 6 until 354:
    ((2 + cosd(5t)) * cosd(2t), (2 + cosd(5t)) * sind(2t)) scaled 0.52cm
    .. endfor cycle);
  addStrandToKnot (cinquefoil) (cord, 0.22cm, "l", "1,-1,1,-1,1,-1,1,-1,1,-1");
  draw knotFromStrands (cinquefoil);
}{Fig.\ 2. The cinquefoil, of five crossings.}
\hfill
\knotcell{
  path cord;
  cord := (for t = 0 step 6 until 354:
    ((2 + cosd(7t)) * cosd(2t), (2 + cosd(7t)) * sind(2t)) scaled 0.52cm
    .. endfor cycle);
  addStrandToKnot (septafoil) (cord, 0.2cm, "l",
    "1,-1,1,-1,1,-1,1,-1,1,-1,1,-1,1,-1");
  draw knotFromStrands (septafoil);
}{Fig.\ 3. The septafoil, of seven crossings.}

\vspace{1.6em}
\small
These are the first three of the family that a cord makes when wound twice
about a ring and turned an odd number of times in going round. Each figure
is one formula sampled into a smooth closed path; the library shades the cord
as a tube, cuts it wherever it passes beneath itself, and lays the strands
one upon another in the order given, so that the eye reads the over and under
of the rope at once.

\vspace{.6em}
\footnotesize
fiziko's own manual calls its knot routine not especially stable. The three
cords here are chosen because their crossings are well separated; a knot
drawn with crossings close together, or with a strand grazing another, may
need its sampling or its width adjusted before the routine finds every
crossing.

\end{document}
